% Дополнения к преамбуле стандартного шаблона pandoc.
% Подключается ко всем целям makefile через -H tex/preamble.tex.

\makeatletter

% --- Десятичная запятая в формулах -----------------------------------------
% icomma делает запятую в математике активной: «2,5» печатается без пробела.
\usepackage{icomma}
% Уточнение правила icomma. В исходном варианте запятая теряет пробел перед
% любым непробельным символом, из-за чего слипаются «f(\nu,T)», «1,\dots»,
% «n=1,2,3,\dots». Здесь запятая считается десятичной, только если сразу за
% ней идёт цифра и за этой группой цифр не следует новая запятая.
\def\icomma@ifdigit{%
  \let\icomma@next\@secondoftwo
  \@tfor\icomma@d:=0123456789\do{%
    \expandafter\ifx\icomma@d\@let@token\let\icomma@next\@firstoftwo\fi}%
  \icomma@next}
\def\sm@rtcomma{%
  \icomma@ifdigit{\let\icomma@digits\@empty\icomma@scan}{\mathcomma}}
\def\icomma@scan#1{%
  \expandafter\def\expandafter\icomma@digits\expandafter{\icomma@digits#1}%
  \futurelet\@let@token\icomma@loop}
\def\icomma@loop{\icomma@ifdigit\icomma@scan\icomma@decide}
\def\icomma@decide{%
  \ifx\@let@token,%
    \let\icomma@next\mathcomma
  \else
    \def\icomma@next{\mathord{\mathcomma}}%
  \fi
  \expandafter\icomma@next\icomma@digits}

% --- Колонтитулы (только для книги) -----------------------------------------
\@ifclassloaded{book}{%
  \usepackage{fancyhdr}
  \pagestyle{fancy}
  \fancyhf{}
  \fancyhead[LE,RO]{\thepage}
  \fancyhead[RE]{\small\nouppercase{\leftmark}}
  \fancyhead[LO]{\small\nouppercase{\rightmark}}
  \renewcommand{\headrulewidth}{0.4pt}
  \setlength{\headheight}{14pt}
  \addtolength{\topmargin}{-2pt}
  % первые страницы глав: без колонтитула, номер страницы внизу
  \fancypagestyle{plain}{\fancyhf{}\fancyfoot[C]{\thepage}%
    \renewcommand{\headrulewidth}{0pt}}
  % марки без заглавных букв
  \renewcommand{\chaptermark}[1]{\markboth{\chaptername\ \thechapter. #1}{}}
  \renewcommand{\sectionmark}[1]{\markright{\thesection. #1}}
  % ненумерованные главы (\chapter*: список литературы) тоже обновляют марки;
  % \@mkboth внутри аргумента (так устроен \tableofcontents) отбрасывается
  \let\orig@schapter\@schapter
  \def\@schapter#1{\orig@schapter{#1}%
    {\let\@mkboth\@gobbletwo\markboth{#1}{#1}}}
}{}

\makeatother
